Ee ass 1 q2 Question Two:
Question Two
i. Using Kirchhoff's voltage law in Figure 1, write an expression for the ideal input base current.
Applying KVL around the base-emitter input loop:
$$V_{BB} - I_B R_B - V_{BE} = 0$$ 
Rearranging the formula to express the base current ($I_B$):
$$\mathbf{I_B = \frac{V_{BB} - V_{BE}}{R_B}}$$ 
ii. Calculate the base current.
Given parameters from Figure 1:
$V_{BB} = 3.3\text{ V}$
$V_{BE} = 0.7\text{ V}$
$R_B = 10\text{ k}\Omega = 10,000\ \Omega$
Substitute the values into the KVL expression:
$$I_B = \frac{3.3\text{ V} - 0.7\text{ V}}{10\text{ k}\Omega} = \frac{2.6\text{ V}}{10\text{ k}\Omega} = \mathbf{0.26\text{ mA}}\ \text{or}\ \mathbf{260\ \mu\text{A}}$$ 
iii. If $V_{CE}$ is $0.5\text{ V}$, calculate the value of the collector current.
Applying KVL around the collector-emitter output loop:
$$V_{CC} - I_C R_C - V_{CE} = 0$$ 
Given parameters from Figure 1:
$V_{CC} = 12\text{ V}$
$V_{CE} = 0.5\text{ V}$
$R_C = 1\text{ k}\Omega = 1,000\ \Omega$
Rearranging to solve for $I_C$:
$$I_C = \frac{V_{CC} - V_{CE}}{R_C}$$ 
$$I_C = \frac{12\text{ V} - 0.5\text{ V}}{1\text{ k}\Omega} = \frac{11.5\text{ V}}{1\text{ k}\Omega} = \mathbf{11.5\text{ mA}}$$ 
iv. A BJT is called a current controlled device. Why is it called so?
A BJT is a current-controlled device because its main output parameter, the collector current ($I_C$), is directly determined and regulated by its input parameter, the base current ($I_B$), according to the fundamental relationship:
$$I_C = \beta \cdot I_B$$ 
(where $\beta$ is the common-emitter current gain factor of the transistor).
v. Find the maximum and minimum values of where this transistor operates.
The operational extremes of the transistor are defined by the Cut-off state and the Saturation state:
Minimum operational limits (Cut-off region):
Occurs when the transistor is fully turned off ($I_B = 0$).
$\mathbf{I_{C(\text{min})} = 0\text{ mA}}$
$\mathbf{V_{CE(\text{max})} = V_{CC} = 12\text{ V}}$
Maximum operational limits (Saturation region):
Occurs when the transistor is fully turned on and conducting maximally. Using the minimum practical saturation voltage given in the text graph ($V_{CE(\text{sat})} \approx 0.5\text{ V}$):
$\mathbf{I_{C(\text{max})} = 11.5\text{ mA}}$ (as calculated in part iii)
$\mathbf{V_{CE(\text{min})} = 0.5\text{ V}}$ (or ideally $0.2\text{ V}$ depending on textbook convention, but matching the provided loop context).
